\documentclass[11pt]{article}
\usepackage[margin=1in]{geometry}
\usepackage{amsmath,amssymb,amsthm,mathtools,mathrsfs,booktabs,array}
\usepackage{microtype,xurl}
\usepackage{xcolor}
\usepackage[colorlinks=true,linkcolor=blue!55!black,citecolor=blue!55!black,urlcolor=blue!55!black]{hyperref}

\newtheorem{theorem}{Theorem}[section]
\newtheorem{proposition}[theorem]{Proposition}
\newtheorem{corollary}[theorem]{Corollary}
\theoremstyle{remark}
\newtheorem{remark}[theorem]{Remark}
\newcommand{\R}{\mathbb R}
\newcommand{\C}{\mathbb C}
\newcommand{\cT}{\mathcal T}
\newcommand{\cQ}{\mathcal Q}
\newcommand{\cO}{\mathcal O}
\newcommand{\rank}{\operatorname{rank}}
\newcommand{\Sol}{\operatorname{Sol}}
\newcommand{\Diff}{\operatorname{Diff}}
\newcommand{\vol}{\operatorname{vol}}
\newcommand{\cert}[1]{\nolinkurl{#1}}

\title{Einstein--Maxwell Waves inside Weyl--Maxwell Gravity\\
on a Compact Product:\\
Exact Linear Phase-Space Inclusion and the Generic Extra Branches}
\author{GPT-5.6.sol\thanks{OpenAI model and principal programme author.}
\and Asger Alstrup Palm\thanks{Programme orchestrator and corresponding human contact; Honorary Professor, DTU Compute, Technical University of Denmark. \texttt{asger@area9.dk}.}}
\date{Working draft, July 2026}

\begin{document}
\maketitle

\begin{abstract}
We locate ordinary Einstein--Maxwell gravitational waves inside pure
Weyl--Maxwell gravity on the compact product
$\R_t\times S^1_L\times S^2$. The background is a common exact solution of
both theories: the sphere has unit curvature, the Maxwell field is a unit
parallel magnetic flux, and the action normalizations are
$\kappa=1$, $\Lambda=1/2$, and $\alpha_B=3$. On the fixed compact $U(1)$
bundle we classify the complete standard linearized Einstein--Maxwell
harmonic phase space, including all axial and polar radiative modes, the
physical exceptional $\ell=1$ quotient, the six-dimensional homogeneous
generalized block, and the three axial twist pairs.

The identity tangent map into Weyl--Maxwell theory is injective and its
Lee--Wald pullback is nondegenerate on every standard block. It is not a
symplectic inclusion. For $\ell\geq2$ the target form is the
Einstein--Maxwell form composed with the parity-independent spectral
polynomial
\[
 p_\lambda(M)=1+\frac32(M-\lambda),
\]
whose two normal-mode weights are $1\pm\frac32\sqrt{2\lambda}$. Thus each
real spatial harmonic has two positive and two negative relative branch
coefficients. The physical $\ell=1$
block instead acquires the factor $4$, the homogeneous block a nontrivial
unipotent shear, and every axial twist pair the factor $-2$.

We then solve the complete generic axial Weyl--Maxwell target. Modulo the
Einstein image, it contains two additional cyclic summands on
$\omega^2-k^2=\lambda-2/3$. Their direct four-dimensional Lee--Wald Gram
matrix is nondegenerate and positive in the declared positive-frequency
current convention. The Einstein and extra modules are mutually orthogonal,
and the normalized positive-frequency current has axial inertia $(3,1)$.
Inverting the extra Gram matrix gives two conserved spectral coefficient
extractors
which annihilate the Einstein image.
An independent literal four-dimensional polar-current calculation certifies
the corresponding generic polar extra block on every declared
$\ell\geq2$ compact-momentum fibre. The generic polar extra block has inertia
$(2,0)$, while the complete generic polar block has inertia $(3,1)$. These
are separate parity calculations: we do not identify axial and polar representatives.

These are classical \textsc{local-algebraic}/\textsc{reduced-mode}
statements before any background-stabilizer/moment-map reduction. They prove
neither a
positive-frequency Hilbert space nor a quantum ghost theorem. They also do
not construct asymptotically flat scattering. The ordinary waves are present
and nonnull at this stage. The absolute $SO(4,2)$ one-particle cohomology of
the conformally flat vacuum cylinder belongs to a different background and
cannot be imported as a quotient theorem for this flux product.
\end{abstract}

\section*{AI-use disclosure and computational accountability}
OpenAI GPT-5 was used substantively for mathematical synthesis, symbolic-code
generation and debugging, literature organization, claim auditing, and
manuscript drafting. The human author posed and refined the scientific
questions, directed the repository workflow, and selected and reviewed the
claims and derivations in this draft. The human author retains responsibility
for the final text, and submission remains conditional on a documented final
human review. Every algebraic result used below is reproduced by deterministic
exact scripts and checked by a separately implemented deterministic verifier
that does not import the generator code; Section~\ref{sec:repro} gives content
hashes and commands. No floating-point calculation is used in a
rank, inertia, quotient, or on-shell identity. This disclosure follows the
current APS requirement that substantive AI use be disclosed while the AI
program is not listed as an author \cite{APSAI}.

\tableofcontents

\section{The question and the answer}\label{sec:intro}

An Einstein--Maxwell solution need not solve Weyl--Maxwell theory, and even
when it does, its tangent solutions need not inherit the Einstein--Maxwell
phase-space form. We study a magnetically supported compactification of the
Pleba\'nski--Hacyan direct-product electrovacuum
\cite{PlebanskiHacyan,KrtousEtAl,Ortaggio}. It lies on the intersection of
the two theories, and we determine both its induced Einstein image and the
generic additional target in each parity. The relevant question is not whether one can fix
Weyl gauge to obtain Einstein gravity---one cannot do so in general
\cite{Bach,Maldacena,DeserJoungWaldron}---but instead:
\begin{quote}
On a common Einstein--Maxwell/Weyl--Maxwell background, which ordinary wave
solutions embed, what phase-space form do they inherit, and which additional
Weyl--Maxwell solutions remain?
\end{quote}

To our knowledge, this paper answers that question for the first time as one
combined phase-space/current/module calculation: completely for the standard harmonic
Einstein--Maxwell tangent and for the independently classified generic axial
and polar parts of the additional target on a fixed compact magnetic bundle.
The answer has three logically
separate layers:
\begin{enumerate}
\item every standard Einstein--Maxwell tangent class injects as a
Weyl--Maxwell solution class;
\item the target Lee--Wald form restricts nondegenerately to that image, but
does not equal the source form under the identity map; and
\item each generic parity sector of the Weyl--Maxwell target contains two
further nonradical solution directions which are not Einstein representatives
or local gauge; the axial and polar representatives are not identified.
\end{enumerate}

Let $\cT_{\rm EM}^{\rm std}$ denote the complete fixed-bundle standard
Einstein--Maxwell harmonic tangent after local $\Diff\times U(1)$ reduction,
and let $\cT_{\rm WM}^{\rm ax}$ denote the generic axial Weyl--Maxwell tangent
after local $\Diff\times\text{Weyl}\times U(1)$ reduction. The central result
can be stated schematically as
\begin{equation}\label{eq:central-summary}
 \boxed{\begin{gathered}
 \iota:\cT_{\rm EM}^{\rm std}\hookrightarrow\cT_{\rm WM},
 \qquad \ker(\iota^*\Omega_{\rm WM})=0,\\
 \iota^*\Omega_{\rm WM}\neq\Omega_{\rm EM},\\
 \left(\cT_{\rm WM}^{\rm ax}/\iota(\cT_{\rm EM}^{\rm ax})\right)_{\ell,n}
 \cong \bigl(K_{\ell,n}[\omega]/(p)\bigr)^2,
 \qquad p=\omega^2-k^2-\lambda+\frac23 .
 \end{gathered}}
\end{equation}
Here
\begin{equation}
 K_{\ell,n}:=\mathbb Q(k_n)\subset\mathbb R,
 \qquad k_n=\frac{2\pi n}{L},
\end{equation}
with $\lambda=\ell(\ell+1)$ already specialized; equivalently one may use
any characteristic-zero field containing these specialized coefficients.
The module statement is algebraic over $K_{\ell,n}[\omega]$, whereas the
current inertia below is a statement on the real positive-frequency solution
space after a real root has been chosen. The last line holds for every
$\ell\geq2$ and every
$n\in\mathbb Z$, including $n=0$; it is a canonical fibrewise module
quotient, not a choice of complement. Section~\ref{sec:extra} proves this
without inverting $k$, $\omega$, or either shell polynomial.

\subsection{What is and is not frozen}
The standard Einstein--Maxwell image is complete in both axial and polar
parities and includes all exceptional/global blocks. The additional generic
$\ell\geq2$ Weyl--Maxwell branch is classified independently in the axial and
polar sectors on the local-gauge-reduced module before the final residual
quotient. Each extra parity block is nonradical with inertia $(2,0)$, and each
complete generic parity block has inertia $(3,1)$. This parity-complete
current classification does not identify axial and polar representatives.
Exceptional extra sectors and every background-stabilizer/moment-map or final
residual descent remain open.

The dependency labels are
\begin{center}
\textsc{local-algebraic}\quad and\quad \textsc{reduced-mode}.
\end{center}
No result below carries a \textsc{lorentzian-causal} dependency label. The
compact background is Lorentzian and the currents are conserved in real
time, but no full Green-hyperbolic BV complex, Hadamard state, null-infinity
phase space, or scattering construction is used.

\section{The two theories and their common compact background}
\label{sec:background}

We use signature $(-,+,+,+)$ and the actions
\begin{align}
 S_{\rm EM}[g,A]
 &=\int_M\!\sqrt{-g}\left[
 \frac{R-2\Lambda}{2\kappa}-\frac14F_{ab}F^{ab}\right]d^4x,
 \label{eq:EM-action}\\
S_{\rm WM}[g,A]
&=\int_M\!\sqrt{-g}\left[
\frac{\alpha_B}{8}C_{abcd}C^{abcd}-\frac14F_{ab}F^{ab}\right]d^4x .
\label{eq:WM-action}
\end{align}
Our Bach convention is fixed by the metric variation
\begin{equation}\label{eq:bach-variation}
 \delta\!\int_M\sqrt{-g}\,C_{abcd}C^{abcd}\,d^4x
 =4\int_M\sqrt{-g}\,B_{ab}\,\delta g^{ab}\,d^4x
 +\int_{\partial M}\theta_C(g;\delta g),
\end{equation}
whereas
$\delta S_{\rm Maxwell}=-\frac12\int\sqrt{-g}\,T_{ab}\delta g^{ab}$
up to its boundary potential. Thus the bulk variation of
\eqref{eq:WM-action} is
$\frac12\int\sqrt{-g}(\alpha_BB_{ab}-T_{ab})\delta g^{ab}$, which fixes both
the factor and sign in the field equation below. The bulk equations are
therefore
\begin{align}
 G_{ab}+\Lambda g_{ab}&=\kappa T_{ab},
 &\nabla_aF^{ab}&=0,\label{eq:EM-eom}\\
 \alpha_B B_{ab}&=T_{ab},
 &\nabla_aF^{ab}&=0,\label{eq:WM-eom}
\end{align}
with $dF=0$ and
\begin{equation}
 T_{ab}=F_{ac}F_b{}^c-\frac14g_{ab}F_{cd}F^{cd}.
\end{equation}

Consider a product $M_2(k_1)\times\Sigma_2(k_2)$ of oriented constant
curvature surfaces and an aligned parallel field
\begin{equation}
 F=E\,\vol(M_2)+P\,\vol(\Sigma_2).
\end{equation}
The same metric and field solve both theories on the incidence branch
\begin{equation}\label{eq:incidence}
 \Lambda=\frac{k_1+k_2}{2},\qquad
 E^2+P^2=\frac{k_2-k_1}{\kappa},\qquad
 \alpha_B\kappa(k_1+k_2)=3.
\end{equation}
This is an intersection of solution loci, not a gauge equivalence.

We specialize to the positive flat branch, locally the magnetically supported
Pleba\'nski--Hacyan universe and globally its spatial circle compactification,
\begin{equation}\label{eq:fixture}
 M=\R_t\times S^1_L\times S^2,\qquad
 (k_1,k_2)=(0,1),\qquad
 (\kappa,\Lambda,\alpha_B,E,P)=(1,\tfrac12,3,0,1).
\end{equation}
The spatial Cauchy surface is the closed manifold $S^1\times S^2$. With
minimal charge $q_{\min}=1$, flux quantization gives Chern number $N=2$.
Throughout the paper the bundle $P_N$ is fixed. Uniform magnetic variation
would change the topological sector and is not an allowed tangent.
The local direct-product background belongs to the family studied in the
Einstein--Maxwell literature \cite{PlebanskiHacyan,KrtousEtAl,Ortaggio} and
in earlier and recent Bach/Weyl--Maxwell product analyses
\cite{FiedlerSchimming,JizbaLeheckova}. The novelty claimed here is not the
existence of the product or of selected exact waves. It is the complete
fixed-bundle standard harmonic phase space, its exact relative Lee--Wald
endomorphism, and the nonradical generic axial quotient.

\subsection{Connections, bundle automorphisms, and global zero modes}
Because $c_1(P_N)\neq0$, the background potential $\bar A$ is a connection on
$P_N$, not a global one-form. For a second connection $A$ on the same bundle,
however,
\begin{equation}
 a=A-\bar A\in\Omega^1(M)
\end{equation}
is a global adjoint-valued one-form. A bundle-covariant infinitesimal lift of
a diffeomorphism and a vertical gauge transformation acts by
\begin{equation}\label{eq:global-gauge}
 \delta_{\xi,\chi}h=\mathcal L_\xi\bar g,\qquad
 \delta_{\xi,\chi}a=\iota_\xi\bar F+d\chi .
\end{equation}
Locally this is the familiar
$\mathcal L_\xi\bar A+d\chi_{\rm local}$ after absorbing
$\iota_\xi\bar A$ into the vertical parameter. Equation
\eqref{eq:global-gauge}, rather than a globally defined $\bar A$, is used when
subtracting the common $\Diff\ltimes\mathcal G_{U(1)}$ action in the
injectivity argument below.

At the infinitesimal identity-component level the constant circle component
$W_x$ is a real tangent coordinate. Globally, if $x$ has period $L$, a large
$U(1)$ transformation shifts
\begin{equation}
 W_x\sim W_x+\frac{2\pi m}{q_{\min}L},\qquad m\in\mathbb Z.
\end{equation}
Thus the finite holonomy is a circle coordinate, while $Q_e$ is its local
cotangent partner. Similarly, the three constant axial twists describe the
tangent to an $SO(3)$ isometry holonomy around $S^1$. Large diffeomorphisms
and conjugation can identify finite points; our theorem concerns the
unreduced local tangent before that global moduli quotient.

\subsection{Why the complete linear inclusion holds}
To remove a convention ambiguity, let
\begin{align}
 \mathsf S_{abcd}={}&-\nabla_aF_{ce}\nabla_bF_d{}^e
 -\nabla_aF_{de}\nabla_bF_c{}^e
 +g_{ab}\nabla_fF_{ce}\nabla^fF_d{}^e\nonumber\\
 &+\frac12g_{cd}\nabla_aF_{ef}\nabla_bF^{ef}
 -\frac14g_{ab}g_{cd}\nabla_eF_{fh}\nabla^eF^{fh},\nonumber\\
 H_{abcd}:={}&\frac12(\mathsf S_{abcd}+\mathsf S_{cdab}),qquad
 H_{ab}:=H_{abc}{}^c,qquad C^{\rm Ch}_{ab}:=2\kappa H_{ab}.
 \label{eq:chevreton-definition}
\end{align}
Thus $C^{\rm Ch}$ is explicitly homogeneous quadratic in $\nabla F$; the
factor $2\kappa$ is part of our convention rather than an unstated change of
the Chevreton tensor.
For a source-free four-dimensional Einstein--Maxwell solution with
cosmological constant, the convention-adjusted trace of the Chevreton tensor
obeys the exact on-shell identity
\begin{equation}\label{eq:chevreton-identity}
 B_{ab}-\frac{2\kappa\Lambda}{3}T_{ab}
 =C^{\rm Ch}_{ab}.
\end{equation}
This is Eq.~(66) of Bergqvist and Eriksson translated to
\eqref{eq:bach-variation}, with the nonzero cosmological term retained
\cite{BergqvistEriksson,BergqvistErikssonSenovilla}.

For a fixed-bundle connection tangent $a\in\Omega^1(M)$, write $f=da$; then
$df=0$ holds identically. Thus the field-strength notation in the following
proposition is the curvature of the potential-based tangent used throughout
the phase-space quotient.

An on-shell identity cannot be differentiated only along actual nonlinear
solution families if one wishes to cover arbitrary Jacobi fields. The
following formal bridge supplies the missing step.

\begin{proposition}[Dual-number formal linearization]
\label{prop:dual-number-bridge}
Let $(h,f)$ be any smooth fixed-bundle solution of the complete linearized
Einstein, Maxwell, and Bianchi equations at $(\bar g,\bar F)$. Then
\begin{equation}\label{eq:formal-chevreton-linearization}
 D\!\left(B_{ab}-\frac{2\kappa\Lambda}{3}T_{ab}
 -C^{\rm Ch}_{ab}\right)_{(\bar g,\bar F)}[h,f]=0.
\end{equation}
No assumption that $(h,f)$ integrates to a real one-parameter family is
required.
\end{proposition}

\begin{proof}
Work in the dual-number algebra
$\mathbb D=\mathbb R[\varepsilon]/(\varepsilon^2)$ and set
$(g_\varepsilon,F_\varepsilon)=(\bar g+\varepsilon h,
\bar F+\varepsilon f)$. The metric is invertible over $\mathbb D$, with
\begin{equation}
 g_\varepsilon^{-1}=\bar g^{-1}
 -\varepsilon\bar g^{-1}h\bar g^{-1}.
\end{equation}
Because the zeroth-order fields solve the nonlinear equations and $(h,f)$
solves their linearizations, the full Einstein, Maxwell, and Bianchi
residuals of $(g_\varepsilon,F_\varepsilon)$ vanish in $\mathbb D$.

Now repeat the derivation of \eqref{eq:chevreton-identity} over
$\mathbb D$. Its first step is the Leibniz expansion of $H_{ab}$; the Maxwell
wave identity then follows by differentiating $dF=0$ and
$\nabla^aF_{ab}=0$ and commuting covariant derivatives. The Ricci terms are
replaced using the Einstein equation, and the remaining curvature derivatives
are converted by the differential Bianchi identity and the geometric
definition of $B_{ab}$. These operations are natural tensor algebra,
covariant differentiation, and curvature commutation, so they remain valid
over $\mathbb D$. Hence \eqref{eq:chevreton-identity} holds modulo
$\varepsilon^2$ for $(g_\varepsilon,F_\varepsilon)$; its
$\varepsilon$ coefficient is \eqref{eq:formal-chevreton-linearization}.
This repeats the on-shell derivation at first formal order rather than
assuming nonlinear integrability.
\end{proof}

\begin{remark}[Functorial ingredients of the formal bridge]
For clarity, the derivation repeated over $\mathbb D$ uses precisely:
\begin{enumerate}
\item the Maxwell wave identity derived from $dF=0$ and
$\nabla^aF_{ab}=0$;
\item curvature commutators for covariant derivatives;
\item the Einstein equation to replace Ricci curvature;
\item the differential Bianchi identity; and
\item the geometric definition of the Bach tensor.
\end{enumerate}
Each step is a natural differential-polynomial identity in the metric,
curvature, field strength, and their covariant derivatives, and therefore
extends functorially from $\mathbb R$ to the dual-number algebra. No curve of
exact real solutions is introduced.
\end{remark}

Consequently, on the formal Einstein--Maxwell shell,
\begin{equation}\label{eq:wm-defect}
 \alpha_BB_{ab}-T_{ab}
 =\alpha_B C^{\rm Ch}_{ab}
 +\left(\frac{2\alpha_B\kappa\Lambda}{3}-1\right)T_{ab}.
\end{equation}
The incidence relation gives $2\alpha_B\kappa\Lambda/3=1$. On the product
fixture $\bar\nabla\bar F=0$, while $C^{\rm Ch}$ is homogeneous quadratic in
$\nabla F$. Therefore
\begin{equation}
 C^{\rm Ch}(\bar g,\bar F)=0,\qquad
 \delta C^{\rm Ch}_{(\bar g,\bar F)}[h,a]=0,
\end{equation}
and linearizing \eqref{eq:wm-defect} along any solution of the complete
linearized Einstein--Maxwell equations gives
$\alpha_B\delta B_{ab}-\delta T_{ab}=0$. The Maxwell equation is identical
in the two theories, so the same pair $(h,a)$ solves the complete linearized
Weyl--Maxwell equations.

Proposition~\ref{prop:dual-number-bridge} covers all formal linearized
solutions, including non-integrable ones. It is nevertheless weaker than an
explicit off-shell factorization by global Euler--Lagrange row operators or
an off-shell BV chain map. The Chevreton structure also predicts its own
limitation: the first possible Chevreton failure is quadratic, so no
nonlinear closure follows from the linear theorem.

Let
\begin{equation}
 \iota[(h,a)]_{\Diff\times U(1)}
 =[(h,a)]_{\Diff\times\mathrm{Weyl}\times U(1)}.
\end{equation}
The map is well defined on the declared linear solution quotients. It is also
injective. Indeed, if an Einstein class becomes target gauge, then after
subtracting its common diffeomorphism and $U(1)$ part it has the form
$(h_{ab},a_a)=(2\sigma\bar g_{ab},0)$. The linearized Einstein--Maxwell rows
imply
\begin{equation}\label{eq:weyl-injectivity}
 -3\Delta_{S^2}\sigma+2\sigma=0.
\end{equation}
On $Y_{\ell m}$ this is multiplication by $3\ell(\ell+1)+2$, hence smoothness
forces $\sigma=0$.

\begin{theorem}[Fixed-bundle tangent inclusion]\label{thm:inclusion}
On the fixed background \eqref{eq:fixture}, the identity tangent map sends
every complete linearized Einstein--Maxwell solution to a linearized
Weyl--Maxwell solution and descends injectively through the respective local
gauge quotients.
\end{theorem}

\begin{remark}
Theorem~\ref{thm:inclusion} does not say that the two theories are gauge
equivalent. In particular, a generic Bach-flat solution need not be
conformally Einstein, and Section~\ref{sec:extra} constructs additional
target directions explicitly.
\end{remark}

\section{The complete standard Einstein--Maxwell harmonic phase space}
\label{sec:source}

Write
\begin{equation}
 \lambda=\ell(\ell+1),\qquad k=k_n=\frac{2\pi n}{L},
 \qquad -\Delta_{S^2}Y_{\ell m}=\lambda Y_{\ell m},
\end{equation}
and $N_{\ell m}=\int_{S^2}|Y_{\ell m}|^2d\Omega>0$. A complex coefficient at
$(n,\ell,m)$ is paired by reality with the coefficient at $(-n,\ell,-m)$.
Equivalently one may use real spherical harmonics and cosine/sine circle
modes. We use the latter convention for real mode counts. The construction
is complementary to the standard gauge-invariant Einstein--Maxwell
master-equation framework \cite{KodamaIshibashi}: the compact product, fixed
magnetic bundle, exceptional harmonics, and global Jordan blocks require the
separate reduction given here.

\subsection{Regular radiative blocks}
For $\ell\geq2$, the axial and polar reductions each have two masters. In the
orders $(H,Q)$ and $(K,U)$ their normal-mode matrices are
\begin{equation}\label{eq:master-matrices}
 M_A=\begin{pmatrix}\lambda&2\\ \lambda&\lambda\end{pmatrix},
 \qquad
 M_P=\begin{pmatrix}\lambda&-2\lambda\\ -1&\lambda\end{pmatrix}.
\end{equation}
Both parities have the exact dispersions
\begin{equation}\label{eq:dispersion}
 \omega_\pm^2=k^2+\lambda\pm\sqrt{2\lambda}.
\end{equation}
The source action-normalized Wronskians are controlled, up to the common
positive factor $N_{\ell m}/2$, by
\begin{equation}\label{eq:source-forms}
 G_A=\begin{pmatrix}\lambda^2&2\lambda\\2\lambda&2\lambda\end{pmatrix},
 \qquad
 G_P=\begin{pmatrix}1&-2\\-2&2\lambda\end{pmatrix}.
\end{equation}
They obey $G_AM_A=M_A^TG_A$ and $G_PM_P=M_P^TG_P$. Their determinants are
$2\lambda^2(\lambda-2)$ and $2(\lambda-2)$, so both are positive for
$\lambda\geq6$. This is positivity of the classical source kinetic branch
weights, not yet a one-particle Hilbert norm.

The branch eigenvectors may be taken as
\begin{equation}
 v_{A,\pm}=(1,\ \pm\sqrt{\lambda/2}),\qquad
 v_{P,\pm}=(\mp\sqrt{2\lambda},\ 1).
\end{equation}
These compact parity masters are the exact representation of the ordinary
coupled photon/graviton-like waves. On the magnetically supported product
the metric and Maxwell coefficients mix; helicity $\pm2$ is therefore a
local/high-frequency interpretation, whereas $(n,\ell,m)$, parity, and
branch are the exact global labels.

\subsection{Exceptional and global blocks}
At $\ell=1$ one branch becomes gauge in each parity. The axial source matrix
has kernel $(H,Q)=(1,-1)$ and physical direction $(1,1)$. The polar source
matrix has kernel $(K,U)=(2,1)$ and quotient coordinate
$\Psi_P=U-K/2$. The surviving axial and polar modes obey
$\omega^2=k^2+4$. They are radiative triplets, not the zero-frequency twist.

The complete homogeneous generalized block is six-dimensional. A convenient
representative is
\begin{equation}\label{eq:hom-rep}
 K=a+bt,\qquad C=at^2+\frac b3t^3+c+dt,\qquad A_x=W_x+Q_et.
\end{equation}
Polynomial Jordan partners are retained; imposing boundedness in time would
define a different phase space. After the common factor $2\pi L$ is removed,
the source form in the order $(a,b,c,d,Q_e,W_x)$ is
\begin{equation}\label{eq:hom-source}
 \Omega_{\rm EM}^{(0)}=
 \begin{pmatrix}
 0&-1&0&-1&0&0\\ 1&0&1&0&0&0\\
 0&-1&0&0&0&0\\ 1&0&0&0&0&0\\
 0&0&0&0&0&-2\\ 0&0&0&0&2&0
 \end{pmatrix}.
\end{equation}
It has rank six. The flat holonomy $W_x$ is the canonical partner of the
electric charge $Q_e$ and must not be deleted merely because it is invisible
to the field strength.

For every one of the three real axial $\ell=1$ harmonics, the generalized
twist
\begin{equation}\label{eq:twist-rep}
 h_{xa}=(A_m+B_mt)X_a,\qquad a_x=-(A_m+B_mt)Y_{1m}
\end{equation}
forms a Darboux pair. The would-be gauge generator is proportional to
$x(A_m+B_mt)$ and is not periodic on $S^1$.

For completeness, Table~\ref{tab:source-strata} records every exceptional
intersection. ``Equation reduction'' counts the surviving coefficient
functions after gauge fixing and then the independent second-order masters;
it is more informative here than the rank of one Fourier matrix, whose rank
jumps at the generalized zero-frequency block. Phase dimensions and current
ranks are real and are stated per real spatial harmonic.

\begin{table}[ht]
\centering
\small
\begin{tabular}{@{}llllcc@{}}
\toprule
Stratum & raw coefficients & gauge rank & equation reduction
  & $\dim\cT$ & $\rank\Omega_{\rm EM}$\\
\midrule
$\ell=0$, $k\neq0$ & $5$ & $2$ & $3\to0$ & $0$ & $0$\\
$\ell=0$, $k=0$ generalized & $5+W_x$ & $2$ & $(a,b,c,d,Q_e,W_x)$
  & $6$ & $6$\\
$\ell=1$, $k\neq0$ & $5_{\rm ax}+7_{\rm pol}$ & $2+3$
  & $3+4\to1+1$ & $4$ & $4$\\
$\ell=1$, $k=0$ & $5_{\rm ax}+7_{\rm pol}$ & $2+3$
  & $1+1$ masters $+$ twist & $6$ & $6$\\
$\ell\geq2$, $k=0$ & $6_{\rm ax}+8_{\rm pol}$ & $2+3$
  & $4+5\to2+2$ & $8$ & $8$\\
$\ell\geq2$, $k\neq0$ & $6_{\rm ax}+8_{\rm pol}$ & $2+3$
  & $4+5\to2+2$ & $8$ & $8$\\
\bottomrule
\end{tabular}
\caption{Complete standard source strata before any optional
background-stabilizer reduction.
At $\ell=0$, every nonzero Fourier block is pure gauge and the electric
coefficient vanishes; exact rank witnesses are $k^4/2$ for $k\neq0$ and
$\omega^4/2$ for $k=0$, $\omega\neq0$. The displayed homogeneous row keeps
the polynomial Jordan solutions rather than imposing temporal boundedness.}
\label{tab:source-strata}
\end{table}

\begin{theorem}[Complete source phase space]\label{thm:source-complete}
Before any optional background-stabilizer/moment-map reduction, the
fixed-bundle standard harmonic
Einstein--Maxwell tangent is the symplectic direct sum
\begin{equation}\label{eq:source-decomp}
 \cT_{\rm EM}^{\rm std}
 =\cT_{\rm rad}^{\ell\geq2}\oplus\cT_{\rm phys}^{\ell=1}
 \oplus\cT_{\rm hom}^{\ell=0}
 \oplus\cT_{\rm twist}^{\ell=1,\omega=0}.
\end{equation}
The regular radiative forms, exceptional physical quotient, homogeneous
block, and three twist pairs are all nondegenerate.
\end{theorem}

\section{The relative phase-space endomorphism on the Einstein image}
\label{sec:restriction}

A solution inclusion need not preserve a symplectic structure. We use the
covariant phase-space current in the Lee--Wald convention
\cite{LeeWald,IyerWald} and therefore
keep three objects separate:
\begin{equation}
 \Omega_{\rm EM},\qquad \iota^*\Omega_{\rm WM},\qquad
 \Omega_{\rm WM}\ \hbox{on the full target}.
\end{equation}
The first two live on the same vector space but arise from different actions.
Because $\Omega_{\rm EM}$ is nondegenerate on the source phase space, define
the relative endomorphism $R$ invariantly by
\begin{equation}\label{eq:relative-endomorphism}
 \iota^*\Omega_{\rm WM}(u,v)=\Omega_{\rm EM}(u,Rv).
\end{equation}
The conjugacy class of $R$, with the identity field map and source
normal-mode identification fixed, is the comparison datum. Abstractly, any
two finite-dimensional nondegenerate real symplectic spaces of the same
dimension are symplectomorphic; our claim is therefore not an absolute
obstruction to some unrelated change of phase-space coordinates.

\subsection{Regular radiation: one spectral polynomial}
On both parities and for every $\ell\geq2$,
\begin{equation}\label{eq:spectral-identity}
 \iota^*\Omega_{\rm WM}(u,v)
 =\Omega_{\rm EM}\bigl(u,p_\lambda(M)v\bigr),
 \qquad p_\lambda(x)=1+\frac32(x-\lambda).
\end{equation}
In the axial and polar master coordinates,
\begin{equation}
 p_\lambda(M_A)=
 \begin{pmatrix}1&3\\3\lambda/2&1\end{pmatrix},\qquad
 p_\lambda(M_P)=
 \begin{pmatrix}1&-3\lambda\\-3/2&1\end{pmatrix}.
\end{equation}
On the normal branches its eigenvalues are
\begin{equation}\label{eq:relative-weights}
 r_+=1+\frac32\sqrt{2\lambda}>0,\qquad
 r_-=1-\frac32\sqrt{2\lambda}<0.
\end{equation}
Neither vanishes for $\lambda\geq6$. The axial and polar branch weights are
identical even though their off-shell coefficient matrices differ.

For a real spatial harmonic there are four oscillator blocks: two parities
times two branches. Hence the inertia of the symmetric branch-coefficient
form associated with $R$ is
\begin{equation}\label{eq:radiative-inertia}
 \operatorname{inertia}_{\rm rel}(R)=(2,2).
\end{equation}
This counts the signs of the symmetric coefficients multiplying the standard
real Darboux blocks. It is not a Sylvester signature of a real symplectic
form, which has no such invariant.

\subsection{Exceptional and global restrictions}
The exceptional $\ell=1$ calculation cannot be obtained by blindly setting
$\lambda=2$ in the generic polar matrix: that continuation fails target
gauge descent. Direct four-dimensional currents instead give, in
source-normalized axial and polar quotient coordinates,
\begin{equation}\label{eq:ell1-factor}
 \left.\iota^*\Omega_{\rm WM}\right|_{\ell=1,\rm phys}
 =4\left.\Omega_{\rm EM}\right|_{\ell=1,\rm phys}.
\end{equation}
The literal target current has zero gauge rows and columns, so this is a
quotient identity rather than only a norm comparison.

On the homogeneous block, after the same common factor as in
\eqref{eq:hom-source},
\begin{equation}\label{eq:hom-target}
 \Omega_{\rm WM}^{(0)}=
 \begin{pmatrix}
 0&2&0&-1&0&0\\ -2&0&1&0&0&0\\
 0&-1&0&0&0&0\\ 1&0&0&0&0&0\\
 0&0&0&0&0&-2\\ 0&0&0&0&2&0
 \end{pmatrix}.
\end{equation}
In the matrix convention of \eqref{eq:relative-endomorphism}, one obtains
\begin{equation}\label{eq:unipotent}
 R=I+N,\qquad N^2=0,\qquad\rank N=2.
\end{equation}
Thus the identity is not symplectic, although $S=I+N/2$ has determinant one
and satisfies $S^T\Omega_{\rm EM}^{(0)}S=\Omega_{\rm WM}^{(0)}$.

Finally, each generalized axial twist pair obeys the direct identity
\begin{equation}\label{eq:twist-factor}
 \left.\iota^*\Omega_{\rm WM}\right|_{\rm twist}
 =-2\left.\Omega_{\rm EM}\right|_{\rm twist}.
\end{equation}
It comes from the full $A+Bt$ current and is not a zero-frequency limit of an
oscillatory formula.

\subsection{Orthogonality and the assembled theorem}
Different circle momenta, spherical irreducibles, and parity sectors are
orthogonal by $S^1$ Fourier invariance, $SO(3)$ invariance, and spatial
parity. Distinct radiative branches are orthogonal because $M$ is
$\Omega_{\rm EM}$-self-adjoint and has distinct eigenvalues. The only shared
label not settled by these arguments is the physical axial
$\ell=1,n=0$ oscillator against the twist. Writing $p_{\rm osc}$ for the
arbitrary physical-oscillator amplitude, its literal mixed target current is
proportional to
\begin{equation}
 -2i\pi\omega p_{\rm osc}(\omega^2-4)
 [\omega(A+Bt)-iB]e^{-i\omega t},
\end{equation}
and therefore vanishes on the physical shell $\omega^2=4$ for both Jordan
coordinates.

\begin{theorem}[Relative phase-space endomorphism]
\label{thm:standard-inclusion}
On the fixed compact background \eqref{eq:fixture}, and before the final
residual quotient, the identity map
\begin{equation}
 \iota:\cT_{\rm EM}^{\rm std}\hookrightarrow\cT_{\rm WM}
\end{equation}
has nondegenerate target pullback:
\begin{equation}
 \ker\left(\iota^*\Omega_{\rm WM}
       \big|_{\cT_{\rm EM}^{\rm std}}\right)=0.
\end{equation}
The pullback is the block-orthogonal direct sum of
\eqref{eq:spectral-identity}, \eqref{eq:ell1-factor},
\eqref{eq:hom-target}, and \eqref{eq:twist-factor}. Equivalently, $R$ is
$p_\lambda(M)$ on regular radiation, $4I$ on the physical $\ell=1$ block,
$I+N$ with $N^2=0$ and $\rank N=2$ on the homogeneous block, and $-2I$ on
each twist pair. The identity inclusion does not preserve the two-forms.
\end{theorem}

\begin{corollary}[Restricted linear Hamiltonians]
On the restricted Einstein image, every linear Einstein--Maxwell Hamiltonian
has a unique target-restricted Hamiltonian vector after application of the
inverse blockwise relative operator.
\end{corollary}

\begin{remark}
The corollary is not an embedding into the full target observable algebra.
Extending a Hamiltonian away from the Einstein image and descending it to a
separately declared moment-map-zero quotient are separate questions.
\end{remark}

\section{The complete generic axial extra branch}\label{sec:extra}

We now solve the fourth-order target rather than restricting it to the
Einstein image. For generic $\ell\geq2$, after local gauge contraction the
axial coefficient order is
\begin{equation}
 \Phi=(H_t,H_x,Q_t,Q_x)^T.
\end{equation}
The contraction inverts only the harmless constant $2$: neither $k$ nor
$\omega$ is inverted, so zero momentum is retained.

Put $s=\omega^2-k^2$. The Einstein master polynomial and the extra polynomial
are
\begin{equation}
 q(s)=(s-\lambda)^2-2\lambda,\qquad p(s)=s-\lambda+\frac23.
\end{equation}
The fraction-field calculation over $\mathbb Q(\lambda,k)[\omega]$ is useful
for discovery but, because it inverts $k$, cannot establish the zero-momentum
specialization. We instead use
\begin{equation}\label{eq:physical-ring}
 R_{\rm phys}=\mathbb Q[\lambda,k,
 \lambda^{-1},(\lambda-2)^{-1},(9\lambda-2)^{-1}]
\end{equation}
and regard the reduced Hessian as a matrix over
$R_{\rm phys}[\omega]$. In particular, $k$, $\omega$, $p$, and $q$ are not
units.

There is a $2\times2$ unit minor $-\lambda^2$. Eliminating its corresponding
block uses only the unit matrices $\operatorname{diag}(\lambda,-\lambda)$
and gives a Schur complement $pT$, where, after removing the common unit
$-3/4$,
\begin{align}
 a={}&k^4+2k^2\lambda-k^2\omega^2+\lambda^2-2\lambda,\nonumber\\
 b={}&k\omega(k^2+2\lambda-\omega^2),\nonumber\\
 d={}&k^2\omega^2-\lambda^2+2\lambda\omega^2+2\lambda-\omega^4,
 \qquad T=-\frac34\begin{pmatrix}a&b\\b&d\end{pmatrix}.
 \label{eq:physical-ring-T}
\end{align}
The entries generate the unit ideal because
\begin{multline}\label{eq:physical-ring-bezout}
 \omega^2(k^2+2\lambda-\omega^2)a
 -k\omega(k^2+2\lambda-\omega^2)b
 -\lambda(\lambda-2)d\\
 =\lambda^2(\lambda-2)^2.
\end{multline}
For a matrix $\mathsf K$, let $I_j(\mathsf K)$ denote the ideal generated by
its $j\times j$ minors. It follows that the determinantal ideals of the
reduced Hessian are
\begin{equation}\label{eq:physical-fitting}
 I_1=(1),\qquad I_2=(1),\qquad I_3=(p),\qquad I_4=(p^2q).
\end{equation}
This is the no-$k$-torsion statement required for specialization.

For every physical fibre
$\lambda=\ell(\ell+1)$, $\ell\geq2$, and
$k=2\pi n/L$, $n\in\mathbb Z$, the resulting PID calculation over
$K_{\ell,n}[\omega]$ has Smith factors
\begin{equation}\label{eq:fibre-smith}
 1,\quad1,\quad p,\quad pq.
\end{equation}
Moreover
$\operatorname{Res}_{\omega}(p,q)=4(9\lambda-2)^2/81$ is nonzero on every
physical fibre. The target module and its primary decomposition are therefore
\begin{equation}\label{eq:target-primary-decomposition}
 \frac{K_{\ell,n}[\omega]}{(p)}\oplus
 \frac{K_{\ell,n}[\omega]}{(pq)}
 \cong
 \left(\frac{K_{\ell,n}[\omega]}{(p)}\right)^2
 \oplus\frac{K_{\ell,n}[\omega]}{(q)}.
\end{equation}
The axial Einstein master module is $K_{\ell,n}[\omega]/(q)$. Its inclusion
is a module map, so its image is annihilated by $q$. Since $q$ is a unit on
each $p$-primary summand, that image lies entirely in the $q$-primary target
summand. The quotient inclusion is injective by
Theorem~\ref{thm:inclusion}; both source and target $q$-primary summands have
$K_{\ell,n}$-dimension $\deg_\omega q=4$. The Einstein image therefore equals
the complete $q$-primary summand. Taking the quotient in
\eqref{eq:target-primary-decomposition} gives
\begin{equation}\label{eq:extra-fibre-module}
 \cQ_{{\rm ax},\ell,n}^{\rm extra}
 :=\left(\cT_{\rm WM}^{\rm ax}/
 \iota(\cT_{\rm EM}^{\rm ax})\right)_{\ell,n}
 \cong\bigl(K_{\ell,n}[\omega]/(p)\bigr)^2.
\end{equation}
No global unimodular Smith transformations over the multivariate ring are
claimed. Equation \eqref{eq:extra-fibre-module}, which is the theorem needed
for the compact physical spectrum, includes $n=0$.

Thus there are two independent additional cyclic summands on
\begin{equation}\label{eq:extra-shell}
 \omega_e^2=k^2+\lambda-\frac23.
\end{equation}
Convenient representatives are
\begin{equation}\label{eq:extra-reps}
 e_1=\begin{pmatrix}-k^2-\lambda\\k\omega_e\\\lambda\\0\end{pmatrix},
 \qquad
 e_2=\begin{pmatrix}-k\omega_e\\k^2-\frac23\\0\\\lambda\end{pmatrix}.
\end{equation}
They are not defined by an arbitrary orthogonal complement: they represent a
canonical quotient supported on a different invariant factor.
At $k=0$ the representatives reduce to columns
$(-\lambda,0,\lambda,0)^T$ and
$(0,-2/3,0,\lambda)^T$; their lower $2\times2$ minor is $\lambda^2$, so the
two extra directions remain independent. Two cyclic summands here mean four
dimensions over $K_{\ell,n}$ before choosing a frequency root, two complex
coefficients on the positive-frequency root, and four real phase-space
dimensions per real spatial harmonic.

\subsection{Direct Lee--Wald pairing}
For any stationary positive-frequency block with time dependence
$e^{-i\omega t}$, $\omega>0$, define the normalized Hermitian current form
\begin{equation}\label{eq:hermitian-current}
 h_+(u,v)=\frac{\Omega_{\rm WM}(\bar u,v)}
 {-i\omega L N_{\ell m}},
\end{equation}
with the time orientation and overall action convention fixed by
\eqref{eq:WM-action}. A different positive normalization of the circle
integral does not change its inertia. Reality and antisymmetry of the real
Lee--Wald form give
\begin{equation}
 \overline{h_+(u,v)}
 =\frac{\Omega_{\rm WM}(u,\bar v)}{i\omega L N_{\ell m}}
 =\frac{\Omega_{\rm WM}(\bar v,u)}{-i\omega L N_{\ell m}}
 =h_+(v,u),
\end{equation}
so $h_+$ is Hermitian on a common positive-frequency shell. In the basis
\eqref{eq:extra-reps} its Gram matrix $G_X$ has entries
\begin{align}
 (G_X)_{11}
 &=\frac{\lambda}{6}\bigl[
 9k^2\lambda^2-12k^2\lambda+4k^2
 +9\lambda^3-18\lambda^2\bigr],\label{eq:gx11}\\
 (G_X)_{12}
 &=\frac{k\lambda\omega_e}{6}(3\lambda-2)^2,\label{eq:gx12}\\
 (G_X)_{22}
 &=\frac{\lambda}{18}\bigl[
 27k^2\lambda^2-36k^2\lambda+12k^2
 +36\lambda-8\bigr].\label{eq:gx22}
\end{align}
On the extra shell,
\begin{equation}\label{eq:extra-det}
 \det G_X=\frac13\lambda^4(\lambda-2)(9\lambda-2)>0.
\end{equation}
The first principal minor is also positive for every real $k$ and
$\lambda\geq6$. Hence $\operatorname{inertia}G_X=(2,0)$.

The Einstein and extra primary modules are mutually Lee--Wald orthogonal.
The two Einstein master branches contribute one positive and one negative
coefficient in the same direct-current convention. Therefore:

\begin{theorem}[Axial target decomposition and current inertia]
\label{thm:axial-target}
For every $\ell\geq2$, every allowed compact momentum
$k=2\pi n/L$, $n\in\mathbb Z$, and before the final
residual quotient, the complete generic axial Weyl--Maxwell solution fibre
is the Lee--Wald orthogonal direct sum of the two Einstein master branches and
the two cyclic extra summands \eqref{eq:extra-fibre-module}. It is
nondegenerate. On the positive-frequency eigenspace, the Hermitian current
form \eqref{eq:hermitian-current} has inertia
\begin{equation}
 \operatorname{inertia}\bigl(h_+|_{\mathcal H^{\rm ax}_+}\bigr)=(3,1).
\end{equation}
The extra block has inertia $(2,0)$; the negative target direction lies on
the minus Einstein-image branch.
\end{theorem}

\begin{theorem}[Independent polar direct-current completion]
\label{thm:polar-target}
For every $\ell\geq2$ and every allowed compact momentum, the literal
four-dimensional polar Lee--Wald current is nondegenerate on the independently
reconstructed local-gauge-reduced generic polar solution module before the final residual quotient.
In the same positive-frequency current convention,
the generic polar extra block has inertia $(2,0)$ and the complete generic
polar block has inertia $(3,1)$. The Einstein and extra polar shell blocks are
orthogonal. This theorem and Theorem~\ref{thm:axial-target} do not identify
axial and polar representatives.
\end{theorem}

\begin{remark}[What the sign does not prove]
Theorem~\ref{thm:axial-target} is a sign theorem for a classical conserved
current after a declared positive-frequency normalization. No compatible
complex structure, positive one-particle inner product, BRST physical Hilbert
space, or quantum time evolution has been constructed. Consequently the
inertia $(3,1)$ is not, by itself, a quantum ghost or unitarity theorem.
The negative target coefficient belongs to $\iota^*\Omega_{\rm WM}$, not to
the independent Einstein--Maxwell source action. Familiar higher-derivative
ghost conclusions \cite{Stelle} require additional quantum structure not
constructed here.
\end{remark}

\subsection{A spectral coefficient extractor for the extra modes}
Nondegeneracy makes the extra branch measurable within the declared reduced
solution block. For an on-shell axial solution $\Phi$, define
\begin{equation}\label{eq:detector}
 \cO_X(\Phi)
 =G_X^{-1}\frac{\Omega_{\rm WM}(\overline e,\Phi)}
 {-i\omega_eL N_{\ell m}}\in\C^2,
\end{equation}
where $e=(e_1,e_2)^T$. The local Green identity gives slice independence, and
direct substitution gives
\begin{equation}
 \cO_X(a_1e_1+a_2e_2)=(a_1,a_2)^T,\qquad \cO_X\circ\iota=0
\end{equation}
on the complete generic axial Einstein image. Thus the additional solutions
are neither a reduced-current radical nor a relabeling of Einstein modes.
This is a conserved mode coefficient functional: it depends on the selected
harmonic, stationary frequency splitting, and representatives. It is not a
local spacetime observable. The extractor has not yet been descended through
any declared background-stabilizer reduction or promoted to a relational,
causal, asymptotic,
Peierls, or quantum observable.

\section{The normalization triangle}\label{sec:normalization}

Higher-derivative currents are sensitive to integration-by-parts and action
normalization conventions. Three independent constructions are therefore
compared. First, the gauge-reduced axial equation operator is formally
self-adjoint. In Fourier convention $\partial_t\mapsto-i\omega$,
$\partial_x\mapsto ik$, its exact quadratic kernel in the order
$(H_t,H_x,Q_t,Q_x)$ is
\begin{equation}\label{eq:kernel}
K=\begin{pmatrix}
-\frac\lambda4 A&-\frac{k\lambda\omega}{4}B&\lambda&0\\
-\frac{k\lambda\omega}{4}B&\frac\lambda4 C&0&-\lambda\\
\lambda&0&k^2+\lambda&k\omega\\
0&-\lambda&k\omega&-\lambda+\omega^2
\end{pmatrix},
\end{equation}
where
\begin{align}
A={}&3k^4+6k^2\lambda-3k^2\omega^2-2k^2
 +3\lambda^2-3\lambda\omega^2-8\lambda+6\omega^2,\\
B={}&3k^2+3\lambda-3\omega^2+4,\\
C={}&3k^2\lambda-3k^2\omega^2-6k^2+3\lambda^2
 -6\lambda\omega^2-8\lambda+3\omega^4+2\omega^2.
\end{align}
It satisfies $K(\omega,k)^\dagger=K(-\omega,-k)^T$ and is the mixed Hessian
of
\begin{equation}
 S^{(2)}_{\rm red}=\frac12\int
 \Phi(-\omega,-k)^TK(\omega,k)\Phi(\omega,k).
\end{equation}

Second, the Green current generated by this Hessian obeys the exact local
Green identity. Third, independent direct variation of the four-dimensional
Weyl--Maxwell action gives a Lee--Wald current whose integral over $S^2$
equals $N_{\ell m}$ times that reduced Green current.

\begin{proposition}[Three-point spectral interpolation]
\label{prop:interpolation}
After the common harmonic norm $N_{\ell m}$ is removed, every entry of the
difference between the direct integrated current and the reduced Green
current is a polynomial in $\lambda$ of degree at most two. It has no
rational denominator or square-root dependence, and $SO(3)$ equivariance
makes it independent of $m$. The direct independent-frequency calculations
at $\lambda=6,12,20$ vanish, so the difference vanishes for every
$\lambda=\ell(\ell+1)$, $\ell\geq2$, and every $m$.
\end{proposition}

The polynomial bound is established before imposing either dispersion
relation: the local curvature current contains at most two contracted sphere
Laplacians after harmonic integration, and all integrations by parts use the
polynomial identity $-\Delta_{S^2}Y=\lambda Y$. Thus no
$\lambda$-dependent normalization is hidden in the interpolation. Covariant
presymplectic potentials admit exact-improvement ambiguities
\cite{HarlowWu}; here $a=A-\bar A$ is global on the fixed bundle, the Cauchy
surface is closed, and integrated exact improvements vanish.

\begin{proposition}[Hessian--Green--Lee--Wald triangle]
\label{prop:triangle}
In the generic axial block, the exact equation operator is the Hessian of the
reconstructed quadratic reduced action; its local Green current equals the
direct integrated four-dimensional Lee--Wald current with the normalization
used in Theorem~\ref{thm:axial-target}.
\end{proposition}

Proposition~\ref{prop:triangle} is sufficient for the linear phase-space and
current-inertia theorems: it fixes the operator, local conservation law, direct
four-dimensional current, and overall action sign independently. A literal
second expansion of the unreduced four-dimensional action density would be a
valuable additional audit, but it is not used or claimed here.

\section{Local modes and background stabilizers}
\label{sec:interpretation}

Three complexes must not be conflated. First, after local
$\Diff\times U(1)$ reduction but before any background-stabilizer/moment-map
reduction, the
metric-containing axial and polar modes \eqref{eq:dispersion}, with their
flux-forced Maxwell dressings, exist and are nonnull for both source and
target currents. These are the exact compact representatives of ordinary
gravitational radiation on this background.

Second, the connected automorphism algebra of this magnetically supported
compactified Pleba\'nski--Hacyan background is only
$\mathbb R H\oplus\mathbb R P_x\oplus\mathfrak{so}(3)$: time translation,
circle translation, and sphere rotations with their bundle compensators.
There is no nonzero Weyl compensator. These generators preserve the Einstein
$q$-primary and extra $p$-primary modules and their Lee--Wald forms, but they
are not universal presymplectic null directions. Consequently no absolute
quotient by them is declared here; one must first construct a common
moment-map/Taub-zero derived sector.

The vanishing absolute one-particle $SO(4,2)$ cohomology of the conformally
flat vacuum cylinder is therefore a statement about a different background
and residual complex. It neither sets the preceding product-background modes
to zero nor determines whether the spectral extractors
\eqref{eq:detector} survive an eventual, separately justified reduction.

Third, $S^1\times S^2$ has no null infinity. An asymptotically flat radiative
phase space changes the allowed transformations and can turn boundary-acting
transformations into charged symmetries. No conclusion about its scattering
states follows from the compact residual quotient studied here.

\paragraph{Einstein is a sector, not a Weyl gauge.}
The exact relation established here is
\begin{equation}
 \Sol_{\rm EM}^{\rm lin,std}
 \subsetneq \Sol_{\rm WM}^{\rm lin,ax+standard}
\end{equation}
on the declared background and harmonic strata. The strictness is witnessed
by \eqref{eq:extra-fibre-module}. Boundary conditions may later select an Einstein
sector, but such selection is additional dynamics plus boundary data, not a
local Weyl gauge fixing.

\section{Open gates and theorem boundary}\label{sec:open}

The following are not prerequisites for the scoped generic-current paper, but they are
required to widen its title-level theorem.
\begin{enumerate}
\item \textbf{Polar ungauged descent and detector.} The exact generic polar
module and its direct four-dimensional Lee--Wald current, orthogonality,
nonradicality, and inertia are now certified. The ungauged BV/Noether lift,
polar coefficient detector, exceptional extra sectors, and final residual
descent remain open.
\item \textbf{Background-stabilizer/moment-map descent.} Complete the
$H$, $P_x$, and $J_i$ descent and quotient only by a subalgebra proved null
on the declared derived sector.  Post-freeze results classify the generic
finite-harmonic $k=0$ cone, one fixed-$|k|$ opposite-momentum cone, and the
exceptional $\ell=1,k=0$ all-$m$ pure no-go.  Multiple $|k|$ fibres,
infinite-mode completion, generalized global mixtures, and the final
stabilizer quotient remain open. Until then physical mode means the
local-gauge-reduced compact phase space, not an absolute residual state space.
\item \textbf{Literal action-density expansion.} This is an independent
normalization audit. The present theorem instead uses the exact
Hessian--Green--direct-Lee--Wald triangle of Proposition~\ref{prop:triangle}.
\item \textbf{Changed-action cyclic repair.} A post-freeze action-layer
certificate now classifies the complete real parity-even local
Einstein--Maxwell action quotient through four derivatives.  Its exact axial
and polar $q$-primary Hessian response does not contain the minimal reduced
source-action repair requested by the relative quantum programme: the
coefficient of $\lambda$ in the axial Maxwell--Maxwell entry and the
coefficient of $\lambda^2$ in the polar Maxwell--Maxwell entry are independent
cokernel witnesses.  Moreover the full $q$--$p$ cross-response system has
rank six, so no nonzero deformation in this bounded action class leaves the
declared $p$ shell unmixed.  This is a scoped four-derivative no-lift theorem,
not a no-go for higher-derivative, nonlocal, or auxiliary-field theories, and
it authorizes no quantum determinant or QME claim.
\item \textbf{Nonlinear extension.} Linear inclusion does not imply that an
Einstein tangent extends to a Weyl--Maxwell family. Fixed-bundle Taub tests
now show a generic pure-extra no-go, finite-harmonic and fixed-$|k|$
mixed-cone extensions, and an exceptional all-$m$ dipole no-go; those belong
to a separate nonlinear charge-fibre paper and successor certificates.
Selected algebraically special Kundt waves on
direct-product electrovacua are known to survive broad higher-order
modifications \cite{KrtousEtAl,Ortaggio}; the open question here is whether
the complete harmonic Einstein tangent, rather than those special exact
families, is nonlinearly closed in this fixed compact theory.
\item \textbf{Causal boundary selection and scattering.} A future theorem
must specify asymptotically flat function spaces, retarded/advanced
complexes, radiative symplectic flux, residual charges, and whether the
Einstein subspace is dynamically closed. Bondi energy and scattering are not
inferred here.  A separate static pure-Weyl black-hole certificate now gives
a regular-horizon family and exact static first law, but supplies no
time-dependent horizon phase space, stability theorem, or implication for
this compact radiative sector.
\end{enumerate}

The strongest present conclusion is therefore:
\begin{quote}
On the fixed compact product, ordinary Einstein--Maxwell waves embed exactly
and nondegenerately in the linear Weyl--Maxwell phase space, but with a
different Hamiltonian form. Weyl--Maxwell theory also has a complete generic
axial branch of genuine, nonradical additional solutions. Neither branch has
yet been classified on a complete moment-map-zero stabilizer quotient or as
an asymptotic
quantum particle space.
\end{quote}

\section{Exact certification and reproducibility}\label{sec:repro}

Every theorem in this manuscript is a synthesis of exact machine-readable
certificates. The companion claim map
\cert{paper/10-compact-einstein-maxwell-weyl-phase-space-claim-map.json}
pins the full SHA-256 hashes and fails closed if an input changes.

\begin{table}[ht]
\centering
\scriptsize
\begin{tabular}{p{0.29\textwidth}p{0.62\textwidth}}
\toprule
Result & Certificate role\\
\midrule
\cert{EINSTEIN_MAXWELL_PRODUCT_INCIDENCE}
& Common exact product background and fixed-flux rational fixture.\\
\cert{EINSTEIN_MAXWELL_CHEVRETON_TANGENT}
& Complete linear on-shell solution inclusion.\\
\cert{EINSTEIN_MAXWELL_CHEVRETON_FORMAL_LINEARIZATION}
& Dual-number bridge from the on-shell derivation to every formal Jacobi
field, without an integrability assumption.\\
\cert{COMPACT_EM_RADIATIVE_SYMPLECTIC_MATCHING}
& Positive source radiative pairing and direct Lee--Wald match.\\
\cert{EINSTEIN_MAXWELL_WEYL_STANDARD_HARMONIC_SYMPLECTIC_INCLUSION}
& Complete standard image, mixed orthogonality, and nondegenerate pullback.\\
\cert{EINSTEIN_MAXWELL_WEYL_AXIAL_OPERATOR}
& Generic target operator and fraction-field invariant factors.\\
\cert{EINSTEIN_MAXWELL_WEYL_AXIAL_PHYSICAL_RING}
& Physical-ring Fitting ideals and every compact-momentum fibre, including
$k=0$.\\
\cert{EINSTEIN_MAXWELL_WEYL_AXIAL_LEE_WALD_COMPLETION}
& Direct four-dimensional current and complete axial current inertia.\\
\cert{EINSTEIN_MAXWELL_WEYL_POLAR_DIRECT_LEE_WALD_COMPLETION_V1}
& Independent direct four-dimensional polar current, orthogonality, and
complete generic polar current inertia.\\
\cert{BRIDGE_PHASE1_EINSTEIN_EXTRA_CONTRIBUTION_V1}
& Publication-independent axial/polar structural crosswalk and claim
boundaries.\\
\cert{EINSTEIN_MAXWELL_WEYL_AXIAL_REDUCED_ACTION_HESSIAN}
& Reduced Hessian and normalization triangle.\\
\cert{EINSTEIN_MAXWELL_WEYL_AXIAL_EXTRA_DETECTOR}
& Conserved extra spectral coefficient extractors before residual descent.\\
\cert{EINSTEIN_MAXWELL_FOUR_DERIVATIVE_ACTION_RESPONSE_V1}
& Complete parity-even four-derivative action quotient, exact axial/polar
$q$-primary and $q$--$p$ Hessian responses, and the reduced repair no-lift
cokernel.\\
\bottomrule
\end{tabular}
\caption{Principal exact inputs. The claim map also pins the exceptional and
mixed-block component certificates.}
\end{table}

The scoped verification commands are listed below; the breakable path
formatting is typographical and each line is one shell command.
\begingroup\footnotesize
\par\noindent\path{python3 bridge/einstein_sector/verify_einstein_maxwell_product_incidence.py}
\par\noindent\path{python3 bridge/einstein_sector/verify_einstein_maxwell_chevreton_tangent.py}
\par\noindent\path{python3 bridge/einstein_sector/verify_einstein_maxwell_chevreton_formal_linearization.py}
\par\noindent\path{python3 bridge/einstein_sector/verify_einstein_maxwell_radiative_symplectic_matching.py}
\par\noindent\path{python3 bridge/einstein_sector/verify_einstein_maxwell_weyl_standard_harmonic_symplectic_inclusion.py}
\par\noindent\path{python3 bridge/einstein_sector/verify_einstein_maxwell_weyl_axial_operator.py}
\par\noindent\path{python3 bridge/einstein_sector/verify_einstein_maxwell_weyl_axial_physical_ring.py}
\par\noindent\path{python3 bridge/einstein_sector/verify_einstein_maxwell_weyl_axial_lee_wald_completion.py}
\par\noindent\path{python3 bridge/einstein_sector/verify_einstein_maxwell_weyl_axial_reduced_action_hessian.py}
\par\noindent\path{python3 bridge/einstein_sector/verify_einstein_maxwell_weyl_axial_extra_detector.py}
\par\noindent\path{python3 bridge/einstein_sector/verify_einstein_maxwell_weyl_polar_direct_lee_wald_completion.py}
\par\noindent\path{python3 paper/generate_10_compact_einstein_maxwell_weyl_phase_space_claim_map.py --check}
\par\noindent\path{python3 paper/verify_10_compact_einstein_maxwell_weyl_phase_space_claim_map.py}
\endgroup
The manuscript itself is checked with two passes of \texttt{pdflatex}.

\section{Conclusion}

The Einstein sector is visible without interpretive guesswork. On a common
compact Einstein--Maxwell/Weyl--Maxwell background, the complete standard
Einstein--Maxwell harmonic tangent injects into the Weyl--Maxwell solution
space and remains nondegenerate under the target Lee--Wald form. The identity
inclusion nevertheless changes that form in every class of blocks: by
branch-dependent weights on regular radiation, a factor four on the physical
$\ell=1$ quotient, a unipotent shear on the homogeneous sector, and a factor
minus two on the twists.

The target is strictly larger. In each generic parity sector its quotient
contains two additional nonradical solution summands. The axial summands are
detected by conserved coefficient extractors; the independently computed
axial and polar positive-frequency currents each have complete inertia
$(3,1)$, with extra-block inertia $(2,0)$.
These statements settle where ordinary waves and the generic extra branches
sit before any optional background-stabilizer/moment-map reduction. They do
not turn a classical sign
into a quantum ghost theorem, and they do not replace the separate causal
problem of selecting an Einstein scattering sector at infinity.

\begin{thebibliography}{99}
\bibitem{Bach}
R.~Bach,
\emph{Zur Weylschen Relativit\"atstheorie und der Weylschen Erweiterung des
Kr\"ummungstensorbegriffs},
Math. Z. \textbf{9} (1921) 110--135,
\href{https://doi.org/10.1007/bf01378338}{doi:10.1007/bf01378338}.

\bibitem{PlebanskiHacyan}
J.~F.~Pleba\'nski and S.~Hacyan,
\emph{Some exceptional electrovac type D metrics with cosmological constant},
J. Math. Phys. \textbf{20} (1979) 1004--1010,
\href{https://doi.org/10.1063/1.524174}{doi:10.1063/1.524174}.

\bibitem{FiedlerSchimming}
B.~Fiedler and R.~Schimming,
\emph{Exact solutions of the Bach field equations of general relativity},
Rep. Math. Phys. \textbf{17} (1980) 15--36,
\href{https://doi.org/10.1016/0034-4877(80)90073-7}%
{doi:10.1016/0034-4877(80)90073-7}.

\bibitem{Chevreton}
M.~Chevreton,
\emph{Sur le tenseur de super\'energie du champ \'electromagn\'etique},
Nuovo Cimento \textbf{34} (1964) 901--913,
\href{https://doi.org/10.1007/bf02812520}{doi:10.1007/bf02812520}.

\bibitem{BergqvistEriksson}
G.~Bergqvist and I.~Eriksson,
\emph{The Chevreton tensor and Einstein--Maxwell spacetimes conformal to
Einstein spaces},
Class. Quantum Grav. \textbf{24} (2007) 3437--3456,
\href{https://arxiv.org/abs/gr-qc/0703073}{arXiv:gr-qc/0703073}.

\bibitem{BergqvistErikssonSenovilla}
G.~Bergqvist, I.~Eriksson, and J.~M.~M.~Senovilla,
\emph{New electromagnetic conservation laws},
Class. Quantum Grav. \textbf{20} (2003) 2663--2668,
\href{https://doi.org/10.1088/0264-9381/20/13/313}%
{doi:10.1088/0264-9381/20/13/313},
\href{https://arxiv.org/abs/gr-qc/0303036}{arXiv:gr-qc/0303036}.

\bibitem{KodamaIshibashi}
H.~Kodama and A.~Ishibashi,
\emph{Master equations for perturbations of generalised static black holes
with charge in higher dimensions},
Prog. Theor. Phys. \textbf{111} (2004) 29--73,
\href{https://arxiv.org/abs/hep-th/0308128}{arXiv:hep-th/0308128}.

\bibitem{KrtousEtAl}
P.~Krtou\v{s}, J.~Podolsk\'y, A.~Zelnikov, and H.~Kadlecov\'a,
\emph{Higher-dimensional Kundt waves and gyratons},
Phys. Rev. D \textbf{86} (2012) 044039,
\href{https://arxiv.org/abs/1201.2813}{arXiv:1201.2813}.

\bibitem{Ortaggio}
M.~Ortaggio,
\emph{Einstein--Maxwell fields as solutions of higher-order theories},
Eur. Phys. J. C \textbf{82} (2022) 1056,
\href{https://arxiv.org/abs/2205.14392}{arXiv:2205.14392}.

\bibitem{JizbaLeheckova}
P.~Jizba and T.~Lehe\v{c}kov\'a,
\emph{Generalized Birkhoff theorems and $2+2$ direct product spacetimes in
Weyl conformal gravity},
Phys. Rev. D \textbf{113} (2026) 044060,
\href{https://doi.org/10.1103/9g8r-lrjk}{doi:10.1103/9g8r-lrjk},
\href{https://arxiv.org/abs/2512.18482}{arXiv:2512.18482}.

\bibitem{LeeWald}
J.~Lee and R.~M.~Wald,
\emph{Local symmetries and constraints},
J. Math. Phys. \textbf{31} (1990) 725--743,
\href{https://doi.org/10.1063/1.528801}{doi:10.1063/1.528801}.

\bibitem{IyerWald}
V.~Iyer and R.~M.~Wald,
\emph{Some properties of Noether charge and a proposal for dynamical black
hole entropy},
Phys. Rev. D \textbf{50} (1994) 846--864,
\href{https://arxiv.org/abs/gr-qc/9403028}{arXiv:gr-qc/9403028}.

\bibitem{HarlowWu}
D.~Harlow and J.-q.~Wu,
\emph{Covariant phase space with boundaries},
JHEP \textbf{10} (2020) 146,
\href{https://arxiv.org/abs/1906.08616}{arXiv:1906.08616}.

\bibitem{Maldacena}
J.~Maldacena,
\emph{Einstein gravity from conformal gravity},
\href{https://arxiv.org/abs/1105.5632}{arXiv:1105.5632}.

\bibitem{DeserJoungWaldron}
S.~Deser, E.~Joung, and A.~Waldron,
\emph{Partial masslessness and conformal gravity},
J. Phys. A \textbf{46} (2013) 214019,
\href{https://arxiv.org/abs/1208.1307}{arXiv:1208.1307}.

\bibitem{Stelle}
K.~S.~Stelle,
\emph{Renormalization of higher-derivative quantum gravity},
Phys. Rev. D \textbf{16} (1977) 953--969,
\href{https://doi.org/10.1103/physrevd.16.953}{doi:10.1103/physrevd.16.953}.

\bibitem{APSAI}
American Physical Society,
\emph{Appropriate use of AI tools}, author policy,
\href{https://journals.aps.org/authors/appropriate-use-ai-tools}%
{journals.aps.org/authors/appropriate-use-ai-tools}, accessed July 2026.
\end{thebibliography}

\end{document}
